\section{Recommended Reading}

\begin{readingbox}
\textbf{Foundation} Read these for orientation and vocabulary discipline. They
help the reader understand why conscious access, wakefulness, and large-scale
coordination matter before confronting the empirical anchors.
\begin{itemize}[nosep]
  \item Stanislas Dehaene, \emph{Consciousness and the Brain}.
  \item Christof Koch, \emph{Consciousness}.
\end{itemize}
\textbf{Modern Research} These are the strong anchors for the core chapter.
Read them as support for signatures, stability architecture, and transition
dynamics rather than as direct proofs of everyday free will.
\begin{itemize}[nosep]
  \item \emph{Dynamical structure-function correlations provide robust and
        generalizable signatures of consciousness in humans}
        (Communications Biology, 2024, strong).
  \item \emph{Cortical connectivity, local dynamics and stability correlates of
        global conscious states} (Communications Biology, 2025, strong).
  \item \emph{Falling asleep follows a predictable bifurcation dynamic}
        (Nature Neuroscience, 2025, strong).
\end{itemize}
\textbf{Reading discipline} Keep empirical consciousness biology separate from
speculative physical theories. Use this chapter to strengthen the manuscript's
state language, not to erase the distinction between strong evidence and
frontier conjecture.
\end{readingbox}
